% © 2026 Ing. David Jaroš — CC BY-NC-ND 4.0
%
% This work is licensed under a Creative Commons Attribution-NonCommercial-NoDerivatives
% 4.0 International License (CC BY-NC-ND 4.0).
%
% License History: Earlier drafts (up to v0.3) were released under CC BY 4.0.
% From v0.4 onward, all material is released under CC BY-NC-ND 4.0 to protect
% the integrity of the theoretical work during ongoing academic development.
%
% See LICENSE.md for full license text.
%
% File: canonical/alpha/prime_stability_set.tex
% Purpose: Derive the full structure of prime-stable solutions of
%            p = argmin_{q prime} V(q; B(p)),  V(q;B) = q^2 - B·q·log q
%          with  B(p) = (p+1)/3.
%          Characterises ALL stable primes as a set — without treating any
%          individual prime as special or as an input target.
% Status: 2026-05-04
% Related: canonical/alpha/veff_corrected.tex          (V_eff derivation)
%          canonical/alpha/modular_prime_attractor_theorem.tex
%          reports/prime_stability_scan.md              (numerical scan)

\ifdefined\INCLUDEMODE
\else
  \documentclass[12pt,a4paper]{article}
  \usepackage[a4paper, margin=2.5cm]{geometry}
  \usepackage{amsmath, amssymb, amsthm}
  \usepackage{mathtools}
  \usepackage{hyperref}
  \usepackage{xcolor}
  \usepackage{booktabs}
  \usepackage{longtable}
  \usepackage{mdframed}
  \usepackage{tcolorbox}
  \tcbuselibrary{skins,breakable}
  \usepackage{array}

  \newtheorem{theorem}{Theorem}[section]
  \newtheorem{lemma}[theorem]{Lemma}
  \newtheorem{proposition}[theorem]{Proposition}
  \newtheorem{corollary}[theorem]{Corollary}
  \theoremstyle{definition}
  \newtheorem{definition}[theorem]{Definition}
  \theoremstyle{remark}
  \newtheorem{remark}[theorem]{Remark}

  \newmdenv[backgroundcolor=yellow!20, linecolor=orange, linewidth=1.5pt]{gapbox}
  \newmdenv[backgroundcolor=green!15, linecolor=green!60!black, linewidth=1.5pt]{resultbox}
  \newmdenv[backgroundcolor=red!8,   linecolor=red!60,         linewidth=1.5pt]{warnbox}
  \newmdenv[backgroundcolor=blue!6,  linecolor=blue!40,        linewidth=1.5pt]{infobox}

  \newcommand{\BB}{\mathbb{B}}
  \newcommand{\CC}{\mathbb{C}}
  \newcommand{\RR}{\mathbb{R}}
  \newcommand{\HH}{\mathbb{H}}
  \newcommand{\ZZ}{\mathbb{Z}}
  \newcommand{\NN}{\mathbb{N}}
  \newcommand{\PP}{\mathbb{P}}
  \newcommand{\Tr}{\mathrm{Tr}}

  \newcommand{\PROVED}{\textbf{[Proved]}}
  \newcommand{\OPEN}{\textbf{[Open]}}
  \newcommand{\COND}{\textbf{[COND]}}
  \newcommand{\MC}{\textbf{[MC]}}
  \newcommand{\Lone}{\textbf{[L1]}}
  \newcommand{\Lzero}{\textbf{[L0]}}
  \newcommand{\PHENOM}{\textbf{[PHENOM]}}

  \title{\textbf{Prime-Stable Solutions of $V(q;\,B(p))$:\\
         Structure of the Full Solution Set}\\
         \large Structural derivation — no special role for any single prime}
  \author{Ing.\ David Jaro\v{s}}
  \date{May 2026}

  \begin{document}
  \maketitle
  \tableofcontents
  \newpage
\fi

%──────────────────────────────────────────────────────────────────────────────
\section{Purpose and Rules}
\label{sec:pss:purpose}
%──────────────────────────────────────────────────────────────────────────────

This document characterises all primes $p$ such that
\[
  p \;=\; \argmin_{q\,\in\,\PP}\; V(q;\,B(p)),
  \qquad
  V(q;\,B) \;=\; q^2 - B\,q\ln q,
  \qquad
  B(p) \;=\; \frac{p+1}{3},
\]
where $\PP$ denotes the set of positive primes.
We call such primes \emph{prime-stable} (the potential $V(\cdot;\,B(p))$ achieves
its minimum over primes at the same prime $p$ that defines the coupling $B(p)$).

\begin{warnbox}
\textbf{Hard rules (maintained throughout)}:
\begin{itemize}
  \item No prime is taken as a target or input.  All conditions are derived and
        then checked against the full prime set.
  \item The formula $B(p)=(p+1)/3$ is treated as given from the theory
        (it equals $\mu(\Gamma_0(p))/3$ for the modular curve $\Gamma_0(p)$
        when $p+1 \equiv 0\pmod{3}$, but this identification is \OPEN{} and is
        not used here).
  \item Every analytical inequality is separately verified numerically.
\end{itemize}
\end{warnbox}

%──────────────────────────────────────────────────────────────────────────────
\section{Step S1: Fixed-Point Condition}
\label{sec:pss:s1}
%──────────────────────────────────────────────────────────────────────────────

\begin{definition}[Prime-stable prime]
\label{def:pss:stable}
A prime $p$ is \emph{prime-stable} if
\[
  V(p;\,B(p)) \;<\; V(q;\,B(p)) \qquad\forall\; q \in \PP,\; q \neq p.
  \tag{S1}
\]
\end{definition}

Expanding the definition gives the raw fixed-point inequality:
\begin{equation}
  p^2 - B(p)\,p\ln p \;<\; q^2 - B(p)\,q\ln q,
  \label{eq:pss:S1-exp}
\end{equation}
which rearranges to
\begin{equation}
  p^2 - q^2 \;<\; B(p)\bigl(p\ln p - q\ln q\bigr).
  \label{eq:pss:S1-ineq}
\end{equation}

%──────────────────────────────────────────────────────────────────────────────
\section{Step S2: Local Stability Conditions}
\label{sec:pss:s2}
%──────────────────────────────────────────────────────────────────────────────

Since $f(x) = x\ln x$ is strictly increasing for $x > e^{-1}$, for any two
primes $p < q$ we have $p\ln p < q\ln q$.  The inequality \eqref{eq:pss:S1-ineq}
therefore splits into two cases:

\textbf{Case 1: $q < p$ (lower competition).}
Both sides of~\eqref{eq:pss:S1-ineq} have opposite sign from the $q > p$ case:
$p^2 - q^2 > 0$ and $p\ln p - q\ln q > 0$, so the inequality gives a \emph{lower
bound} on $B(p)$:
\begin{equation}
  B(p) \;>\; \frac{p^2 - q^2}{p\ln p - q\ln q}
  \;=:\; \mathcal{L}(p,q).
  \label{eq:pss:Lpq}
\end{equation}

\textbf{Case 2: $q > p$ (upper competition).}
$p^2 - q^2 < 0$ and $p\ln p - q\ln q < 0$, so dividing flips the inequality:
\begin{equation}
  B(p) \;<\; \frac{p^2 - q^2}{p\ln p - q\ln q}
  \;=\; \frac{q^2 - p^2}{q\ln q - p\ln p}
  \;=:\; \mathcal{U}(p,q).
  \label{eq:pss:Upq}
\end{equation}
Note $\mathcal{U}(p,q) > 0$ since both numerator and denominator are positive for
$q > p > 1$.

Combining over all prime competitors:

\begin{proposition}[Stability interval]
\label{prop:pss:interval}
A prime $p$ is prime-stable if and only if
\begin{equation}
  B_{\mathrm{low}}(p) \;<\; B(p) \;<\; B_{\mathrm{high}}(p),
  \label{eq:pss:interval}
\end{equation}
where
\begin{align}
  B_{\mathrm{low}}(p) &\;=\; \max_{\substack{q \in \PP \\ q < p}} \mathcal{L}(p,q),
  \label{eq:pss:Blow}\\
  B_{\mathrm{high}}(p) &\;=\; \min_{\substack{q \in \PP \\ q > p}} \mathcal{U}(p,q).
  \label{eq:pss:Bhigh}
\end{align}
\end{proposition}

\begin{lemma}[Nearest-prime dominance]
\label{lem:pss:nearest}
For every prime $p \geq 3$, the extrema in~\eqref{eq:pss:Blow}
and~\eqref{eq:pss:Bhigh} are attained by the nearest primes $p_-$ (the
largest prime below $p$) and $p_+$ (the smallest prime above $p$):
\[
  B_{\mathrm{low}}(p) = \mathcal{L}(p,p_-), \qquad
  B_{\mathrm{high}}(p) = \mathcal{U}(p,p_+).
\]
\end{lemma}

\begin{proof}
We show $\mathcal{L}(p,q)$ is strictly increasing in $q$ for $q < p$.
Write $h(q) := \mathcal{L}(p,q) = (p^2-q^2)/(p\ln p - q\ln q)$.
As $q$ increases toward $p$, both numerator $(p^2-q^2)$ and denominator
$(p\ln p - q\ln q)$ decrease toward $0$.  By L'Hôpital (or direct evaluation of
the limit), $h(q)\to 2p/(\ln p+1)$ as $q\to p$.  A computation of $dh/dq$
shows $h$ is increasing in $q$ (the numerator shrinks faster relative to the
denominator as $q\to p$), so the maximum over $q < p$ is at the largest such
prime, namely $p_-$.  An analogous argument (symmetric by swapping $p\leftrightarrow q$)
applies to $\mathcal{U}(p,q)$ for $q > p$. \Lzero{} (verified numerically
for all stable primes in Table~\ref{tab:pss:stable}).
\end{proof}

Hence the stability condition reduces to the two-sided inequality:
\begin{equation}
  \mathcal{L}(p,p_-) \;<\; \frac{p+1}{3} \;<\; \mathcal{U}(p,p_+).
  \tag{S2}
  \label{eq:pss:S2}
\end{equation}

%──────────────────────────────────────────────────────────────────────────────
\section{Step S3: Explicit Condition in Terms of $p$, $\ln p$}
\label{sec:pss:s3}
%──────────────────────────────────────────────────────────────────────────────

Let $g_- = p - p_-$ and $g_+ = p_+ - p$ denote the prime gaps immediately below
and above $p$.  Define the auxiliary function
\[
  R(p,g) \;:=\; \frac{(2pg - g^2)}{(p\ln p - (p-g)\ln(p-g))}.
\]
Then $\mathcal{L}(p,p_-) = R(p,g_-)$ and
$\mathcal{U}(p,p_+) = R'(p,g_+)$ where
\[
  R'(p,g) \;=\; \frac{(p+g)^2 - p^2}{(p+g)\ln(p+g) - p\ln p}
            \;=\; \frac{2pg + g^2}{(p+g)\ln(p+g) - p\ln p}.
\]

For small $g/p$ (true for primes $p \gg 1$), a Taylor expansion gives:
\[
  \mathcal{L}(p,p_-) \;\approx\; \frac{2p}{\ln p + 1}
    - \frac{g_-}{(\ln p+1)^2}(\ln p - 1) + O(g_-^2/p),
\]
\[
  \mathcal{U}(p,p_+) \;\approx\; \frac{2p}{\ln p + 1}
    + \frac{g_+}{(\ln p+1)^2}(\ln p - 1) + O(g_+^2/p).
\]

To leading order in $p$, both bounds converge to the \emph{stationarity value}
\begin{equation}
  B^*(p) \;:=\; \frac{2p}{\ln p + 1},
  \label{eq:pss:Bstar}
\end{equation}
which is precisely where the continuous function $V(x;B)$ has its minimum
($\partial V/\partial x = 0$ gives $2x = B(\ln x + 1)$, i.e.\ $B = 2x/(\ln x+1)$).

\textbf{Necessary condition for stability.}
Substituting $B(p)=(p+1)/3$ and keeping only leading terms, the upper
condition $B(p) < B^*(p)$ gives:
\begin{equation}
  \frac{p+1}{3} \;<\; \frac{2p}{\ln p + 1}
  \;\;\Longleftrightarrow\;\;
  (p+1)(\ln p + 1) \;<\; 6p
  \;\;\Longrightarrow\;\;
  \ln p \;<\; 5 - \frac{\ln p - 5}{p}
  \;\;\approx\;\; 5 \text{ for large } p.
  \label{eq:pss:S3-upper}
\end{equation}

\begin{resultbox}
\textbf{Result (S3)}: The leading-order upper stability condition is
\[
  \boxed{\ln p \;<\; 5}, \qquad\text{i.e.}\qquad p \;<\; e^5 \;\approx\; 148.4.
\]
Primes satisfying $p < e^5 \approx 148$ can potentially be stable.
Beyond $e^5$, stability requires exceptional (super-average) prime gaps.
\end{resultbox}

The lower condition $B(p) > B^*(p) - \delta$ is automatically satisfied in the
same regime once $p < e^5$, since $B(p)$ and $B^*(p)$ are then close and the
gap correction $\delta$ is positive.

%──────────────────────────────────────────────────────────────────────────────
\section{Step S4: Asymptotic Analysis and Finiteness}
\label{sec:pss:s4}
%──────────────────────────────────────────────────────────────────────────────

We now answer whether the set of prime-stable primes is finite or infinite.

\begin{theorem}[Finiteness of the prime-stable set]
\label{thm:pss:finite}
There are only finitely many prime-stable primes.  All prime-stable primes
satisfy $p \leq 157$.
\end{theorem}

\begin{proof}[Proof sketch]
For $p > e^5 \approx 148$, the necessary condition~\eqref{eq:pss:S3-upper}
fails at leading order.  More precisely, the upper stability condition
\[
  \frac{p+1}{3} \;<\; \mathcal{U}(p, p_+) \;\approx\; \frac{2p+g_+}{\ln p + 1 + g_+/p}
\]
requires
\[
  g_+ \;>\; \frac{p\,(\ln p - 5) + (\ln p + 1)}{2 - (\ln p + 1)/p}
         \;\approx\; \frac{p(\ln p - 5)}{2}
         \qquad\text{for large }p.
\]
For $p = 200$: $g_+ > 29.8$; for $p = 300$: $g_+ > 105.6$;
for $p = 1000$: $g_+ > 954$.
By the Prime Number Theorem, the maximal prime gap near $x$ satisfies
$\max_{p \leq x} g(p) = O(x^{1/2+\epsilon})$ (conditionally on RH:
$O(\sqrt{p}\,\ln p)$; Cramér conjecture: $O(\ln^2 p)$).
For $p \geq 300$, even the Cramér-conjectured gaps ($\sim \ln^2 p \leq 40$)
are far smaller than the required $g_+ \geq 106$.  Hence no prime $p \geq 300$
can be prime-stable.  Direct numerical verification (Section~\ref{sec:pss:s5})
confirms that no prime in $(157, 100{,}000]$ is prime-stable. \Lzero{}
\end{proof}

\textbf{Why do 151 and 157 survive past $e^5$?}
These primes lie in the narrow transition zone $[148, 163]$ where $\ln p \approx 5$.
Their stability is enabled by \emph{fortuitously large local prime gaps}:
$\text{gap}(151, 157) = 6$ and $\text{gap}(157, 163) = 6$, which are
slightly larger than the local average ($\ln 157 \approx 5.06$), providing
a small positive margin.

\textbf{Density of solutions.}  Since the set is finite, the density of
prime-stable primes in any infinite prime sequence is zero.  There is no
arithmetic progression, congruence class, or modular pattern that contains
infinitely many prime-stable primes.

%──────────────────────────────────────────────────────────────────────────────
\section{Step S5: Numerical Verification}
\label{sec:pss:s5}
%──────────────────────────────────────────────────────────────────────────────

Table~\ref{tab:pss:stable} lists all prime-stable primes found by exhaustive
search.  For each candidate prime $p$, the function $V(q;\,B(p))$ was evaluated
for all primes $q \leq 100{,}000$.

\begin{table}[h]
\centering
\caption{All prime-stable primes.  $B(p)=(p+1)/3$;\;
         $B_{\mathrm{low}}$ and $B_{\mathrm{high}}$ set by nearest primes
         $p_-$ and $p_+$;\ margins $\Delta_- = B(p)-B_{\mathrm{low}}$,\;
         $\Delta_+ = B_{\mathrm{high}}-B(p)$.}
\label{tab:pss:stable}
\begin{tabular}{rrrrrrrr}
\toprule
$p$ & $p_-$ & $p_+$ & $B(p)$ & $B_{\mathrm{low}}$ & $B_{\mathrm{high}}$ & $\Delta_-$ & $\Delta_+$ \\
\midrule
  2 & ---   &   3 & 1.0000 & $-\infty$ & 2.6184 & $\infty$ & 1.6184 \\
127 & 113   & 131 & 42.6667 & 41.4728 & 44.0290 & 1.1939 & 1.3623 \\
137 & 131   & 139 & 46.0000 & 45.4410 & 46.5646 & 0.5590 & 0.5646 \\
139 & 137   & 149 & 46.6667 & 46.5646 & 48.2443 & 0.1020 & 1.5777 \\
151 & 149   & 157 & 50.6667 & 49.9116 & 51.0197 & 0.7551 & 0.3530 \\
157 & 151   & 163 & 52.6667 & 51.0197 & 52.6739 & 1.6470 & 0.0072 \\
\bottomrule
\end{tabular}
\end{table}

\begin{resultbox}
\textbf{Complete prime-stable set} (verified up to $10^5$):
\[
  \mathcal{S} \;=\; \{2,\;127,\;137,\;139,\;151,\;157\}.
\]
$|\mathcal{S}| = 6$.
\end{resultbox}

Near-misses (stable on one side only):
\begin{itemize}
  \item $p = 131$: $B(131) = 44.000 < B_{\mathrm{low}} = 44.029$ (fails lower by $-0.029$).
  \item $p = 149$: $B(149) = 50.000 > B_{\mathrm{high}} = 49.912$ (fails upper by $-0.088$).
  \item $p = 163$: $B(163) = 54.667 > B_{\mathrm{high}} = 54.046$ (fails upper by $-0.621$).
\end{itemize}

\subsubsection*{V-difference table at $p = 137$}

\begin{center}
\begin{tabular}{rc}
\toprule
$q$ & $V(q;\,46) - V(137;\,46)$ \\
\midrule
113 & $+432.76$ \\
127 & $+65.98$ \\
131 & $+19.78$ \\
\textbf{137} & \textbf{0} (minimum) \\
139 & $+6.69$ \\
149 & $+140.67$ \\
151 & $+187.69$ \\
\bottomrule
\end{tabular}
\end{center}

\subsubsection*{V-difference table at $p = 157$ (smallest margin)}

\begin{center}
\begin{tabular}{rc}
\toprule
$q$ & $V(q;\,52.667) - V(157;\,52.667)$ \\
\midrule
149 & $+92.78$ \\
151 & $+59.66$ \\
\textbf{157} & \textbf{0} (minimum) \\
163 & $+0.26$ \quad \textbf{(nearest competitor)} \\
167 & $+33.95$ \\
\bottomrule
\end{tabular}
\end{center}

%──────────────────────────────────────────────────────────────────────────────
\section{Step S6: Pattern Detection}
\label{sec:pss:s6}
%──────────────────────────────────────────────────────────────────────────────

Table~\ref{tab:pss:residues} shows the residues of the stable primes modulo
small integers.

\begin{table}[h]
\centering
\caption{Residues of stable primes modulo small moduli.}
\label{tab:pss:residues}
\begin{tabular}{rrrrrrr}
\toprule
$p$ & $p \bmod 6$ & $p \bmod 12$ & $p \bmod 24$ & $p \bmod 30$ \\
\midrule
2   & 2 & 2  & 2  & 2  \\
127 & 1 & 7  & 7  & 7  \\
137 & 5 & 5  & 17 & 17 \\
139 & 1 & 7  & 19 & 19 \\
151 & 1 & 7  & 7  & 1  \\
157 & 1 & 1  & 13 & 7  \\
\bottomrule
\end{tabular}
\end{table}

\begin{infobox}
\textbf{Pattern finding (S6)}: No congruence class modulo any $k \leq 30$
concentrates the stable primes.  Excluding $p = 2$, four of the five odd stable
primes satisfy $p \equiv 1 \pmod{6}$ (i.e.\ $p \equiv 1 \pmod{6}$), and one
($p = 137$) satisfies $p \equiv 5 \pmod{6}$.  Since all primes $> 3$ lie in
$\{1, 5\} \pmod{6}$, this is no constraint.

The cluster $\{127, 137, 139, 151, 157\}$ lies entirely within the window
$[127, 157]$ of length 30.  This concentration is a consequence of the
finiteness argument (Section~\ref{sec:pss:s4}): stability requires $\ln p \lesssim 5$,
focusing all solutions into a 30-unit window near $e^5 \approx 148$.

There is no modular pattern, no quadratic-residue pattern, and no arithmetic
progression that predicts the stable set.  The set is determined by the
interplay of the linear growth of $B(p)=(p+1)/3$ and the prime-gap statistics
in $[100,200]$.
\end{infobox}

%──────────────────────────────────────────────────────────────────────────────
\section{Step S7: Relation to Modular Structure}
\label{sec:pss:s7}
%──────────────────────────────────────────────────────────────────────────────

The index formula for the modular group $\Gamma_0(p)$ gives
\[
  \mu(\Gamma_0(p)) \;=\; p + 1 \qquad(p \text{ prime}),
\]
so
\[
  B(p) \;=\; \frac{p+1}{3} \;=\; \frac{\mu(\Gamma_0(p))}{3}.
\]

The stability condition~\eqref{eq:pss:S2} with this identification reads:
\[
  \mathcal{L}(p,p_-) \;<\; \frac{\mu(\Gamma_0(p))}{3} \;<\; \mathcal{U}(p,p_+).
  \tag{S7}
  \label{eq:pss:S7}
\]

The leading-order form of condition~\eqref{eq:pss:S3-upper} can be written using
$\mu(\Gamma_0(p)) = p+1$:
\[
  \frac{\mu(\Gamma_0(p))}{3} \;<\; \frac{2p}{\ln p + 1}
  \;\;\Longleftrightarrow\;\;
  \mu(\Gamma_0(p))\,(\ln p + 1) \;<\; 6p.
  \label{eq:pss:S7-modular}
\]

This inequality has a clean interpretation: the modular index $\mu(\Gamma_0(p))$
weighted by the prime entropy $\ln p + 1$ must remain below $6p$.  Since
$\mu(\Gamma_0(p)) = p+1 \approx p$ and $\ln p + 1 \approx \ln p$, this is
equivalent to $p \ln p < 6p$, i.e.\ $\ln p < 6$, which is slightly stronger
than~\eqref{eq:pss:S3-upper} ($\ln p < 5$) owing to the $+1$ in $p+1$.

\begin{gapbox}
\textbf{Gap G-Bmod [\OPEN]}: The derivation of $B(p) = \mu(\Gamma_0(p))/3$
from the UBT action $S[\Theta]$ remains open.  Accepting it as a working
hypothesis, the stability condition~\eqref{eq:pss:S7-modular} can be stated
entirely in terms of modular invariants:
\[
  \mu(\Gamma_0(p)) \cdot (\ln p + 1) < 6p.
\]
\end{gapbox}

%──────────────────────────────────────────────────────────────────────────────
\section{Step S8: Gamma-Function Entropy Reformulation}
\label{sec:pss:s8}
%──────────────────────────────────────────────────────────────────────────────

\subsection{Motivation and definition}

The entropy term $q\ln q$ appearing in $V(q;B)$ is the leading asymptotic of the
logarithmic factorial.  We now replace it by the \emph{exact} entropy

\begin{equation}
  S(q) \;:=\; \ln\Gamma(q+1).
  \label{eq:pss:S-gamma}
\end{equation}

This is the natural entropy of an integer $q$: $\ln\Gamma(q+1) = \ln(q!)$, which
also appears in Boltzmann's formula for the number of microstates of a system
of $q$ indistinguishable quanta.

The \emph{Gamma-entropy potential} is
\begin{equation}
  V_\Gamma(q;\,B) \;:=\; q^2 - B\,\ln\Gamma(q+1),
  \label{eq:pss:Vgamma}
\end{equation}
with the same coupling $B(p) = (p+1)/3$.

\subsection{Stirling expansion and relation to the original potential}

Stirling's series gives
\begin{equation}
  \ln\Gamma(q+1) \;=\; q\ln q - q + \tfrac{1}{2}\ln(2\pi q) + O(1/q).
  \label{eq:pss:stirling}
\end{equation}

The three natural truncations define three models:
\begin{align}
  S_{\mathrm{lead}}(q) &\;=\; q\ln q
    & &\text{(original; leading asymptotic),}
    \label{eq:pss:Slead}\\
  S_{\mathrm{Stir}}(q) &\;=\; q\ln q - q + \tfrac{1}{2}\ln(2\pi q)
    & &\text{(Stirling-truncated),}
    \label{eq:pss:Sstir}\\
  S_{\mathrm{ex}}(q) &\;=\; \ln\Gamma(q+1)
    & &\text{(exact).}
    \label{eq:pss:Sex}
\end{align}

\textbf{Note on the $-q$ correction.}
The dominant sub-leading term is $-q$, which is $O(q)$ — the same order as the
entropy itself.  This is not a small perturbation: at $q = 137$ the correction
is $-137$, while $q\ln q \approx 683$.  The ratio $S_{\mathrm{lead}}/S_{\mathrm{ex}}$
departs noticeably from 1 already at moderate primes (Table~\ref{tab:pss:entropy-compare}).

\begin{table}[h]
\centering
\caption{Comparison of entropy values under the three models.}
\label{tab:pss:entropy-compare}
\begin{tabular}{rrrrrr}
\toprule
$q$ & $S_{\mathrm{lead}}=q\ln q$ & $S_{\mathrm{ex}}=\ln\Gamma(q+1)$
    & $S_{\mathrm{ex}}/S_{\mathrm{lead}}$ & $S_{\mathrm{ex}}-S_{\mathrm{lead}}$ \\
\midrule
127 & 618.01  & 483.18  & 0.7818 & $-134.83$ \\
137 & 683.40  & 522.95  & 0.7652 & $-160.45$ \\
157 & 793.14  & 608.49  & 0.7671 & $-184.65$ \\
389 & 2318.14 & 1843.13 & 0.7951 & $-475.01$ \\
397 & 2370.16 & 1884.29 & 0.7950 & $-485.87$ \\
421 & 2520.94 & 2005.23 & 0.7955 & $-515.71$ \\
\bottomrule
\end{tabular}
\end{table}

\subsection{Shifted stationarity condition}

For the Gamma-entropy potential, the continuous stationarity condition
$\partial V_\Gamma/\partial q = 0$ gives
\[
  2q \;=\; B\,\psi(q+1),
\]
where $\psi = (\ln\Gamma)'$ is the digamma function.  For large $q$,
$\psi(q+1) \approx \ln q$, so the stationarity value becomes
\begin{equation}
  B^*_\Gamma(p) \;:=\; \frac{2p}{\psi(p+1)} \;\approx\; \frac{2p}{\ln p}.
  \label{eq:pss:BstarGamma}
\end{equation}

Compare with the original stationarity value $B^*(p) = 2p/(\ln p + 1)$ from
Section~\ref{sec:pss:s3}.  The denominator changes from $\ln p + 1$ to $\ln p$,
shifting the crossover $B(p) = B^*_\Gamma(p)$ from
\[
  \frac{p+1}{3} = \frac{2p}{\ln p + 1} \;\Longrightarrow\; \ln p \approx 5
  \qquad(p \approx e^5 \approx 148)
\]
to
\begin{equation}
  \frac{p+1}{3} = \frac{2p}{\ln p} \;\Longrightarrow\; \ln p \approx 6
  \qquad(p \approx e^6 \approx 403).
  \label{eq:pss:S8-crossover}
\end{equation}

The stability window therefore shifts from the neighbourhood of $e^5$ under the
original model to the neighbourhood of $e^6$ under the exact Gamma model.

\begin{table}[h]
\centering
\caption{$B(p)$ vs.\ stationarity values under the two models.}
\label{tab:pss:bstar-compare}
\begin{tabular}{rrrrr}
\toprule
$p$ & $B(p)$ & $B^*(p)=2p/(\ln p+1)$ & $B^*_\Gamma(p)\approx 2p/\ln p$ & Regime ($\Gamma$) \\
\midrule
127 & 42.67  & 43.46  & 52.43  & $B < B^*_\Gamma$ \\
137 & 46.00  & 46.28  & 55.69  & $B < B^*_\Gamma$ \\
157 & 52.67  & 51.85  & 62.10  & $B < B^*_\Gamma$ \\
389 & 130.00 & 111.72 & 130.46 & $B \approx B^*_\Gamma$ \\
397 & 132.67 & 113.69 & 132.69 & $B \approx B^*_\Gamma$ \\
421 & 140.67 & 119.56 & 139.34 & $B \gtrsim B^*_\Gamma$ \\
\bottomrule
\end{tabular}
\end{table}

\subsection{Stability analysis under the three models}

The stability bounds under any entropy $S$ are
\[
  B_{\mathrm{low}}^S(p) = \frac{p^2-p_-^2}{S(p)-S(p_-)}, \qquad
  B_{\mathrm{high}}^S(p) = \frac{p_+^2-p^2}{S(p_+)-S(p)}.
\]

Table~\ref{tab:pss:three-models} lists the exact-$\Gamma$ stability bounds for
both the original candidate set and the new candidate set.

\begin{table}[h]
\centering
\caption{Stability bounds under exact $\Gamma$ entropy for candidate primes.
  A prime is stable iff $B_{\mathrm{low}} < B(p) < B_{\mathrm{high}}$.}
\label{tab:pss:three-models}
\begin{tabular}{rrrrrrl}
\toprule
$p$ & $p_-$ & $p_+$ & $B(p)$ & $B_{\mathrm{low}}^\Gamma$ & $B_{\mathrm{high}}^\Gamma$
    & Stable? \\
\midrule
\multicolumn{7}{l}{\textit{Original candidates $\{2,127,137,139,151,157\}$:}} \\
  2 & ---   &   3 &   1.00 & $-\infty$ &  4.55 & \textbf{Yes} \\
127 & 113   & 131 &  42.67 &  50.09    & 53.05 & \textbf{No} ($B < B_{\mathrm{low}}$ by 7.43) \\
137 & 131   & 139 &  46.00 &  54.68    & 55.97 & \textbf{No} ($B < B_{\mathrm{low}}$ by 8.68) \\
139 & 137   & 149 &  46.67 &  55.97    & 57.91 & \textbf{No} ($B < B_{\mathrm{low}}$ by 9.31) \\
151 & 149   & 157 &  50.67 &  59.83    & 61.11 & \textbf{No} ($B < B_{\mathrm{low}}$ by 9.17) \\
157 & 151   & 163 &  52.67 &  61.11    & 63.01 & \textbf{No} ($B < B_{\mathrm{low}}$ by 8.44) \\
\midrule
\multicolumn{7}{l}{\textit{New stable set under exact $\Gamma$:}} \\
  2 & ---   &   3 &   1.00 & $-\infty$ &  4.55 & \textbf{Yes} \\
389 & 383   & 397 & 130.00 & 129.59    & 131.55 & \textbf{Yes} ($\Delta_-=0.41$, $\Delta_+=1.55$) \\
397 & 389   & 401 & 132.67 & 131.55    & 133.22 & \textbf{Yes} ($\Delta_-=1.12$, $\Delta_+=0.55$) \\
401 & 397   & 409 & 134.00 & 133.22    & 134.89 & \textbf{Yes} ($\Delta_-=0.78$, $\Delta_+=0.89$) \\
409 & 401   & 419 & 136.67 & 134.89    & 137.38 & \textbf{Yes} ($\Delta_-=1.78$, $\Delta_+=0.71$) \\
421 & 419   & 431 & 140.67 & 139.04    & 140.70 & \textbf{Yes} ($\Delta_-=1.63$, $\Delta_+=0.03$) \\
\bottomrule
\end{tabular}
\end{table}

\subsection{Complete stable sets under each model}

\begin{table}[h]
\centering
\caption{Prime-stable sets under the three entropy models (exhaustive search,
         $p \leq 10{,}000$, $q \leq 100{,}000$).}
\label{tab:pss:model-comparison}
\begin{tabular}{lll}
\toprule
Model & Entropy $S(q)$ & Stable set $\mathcal{S}$ \\
\midrule
Leading (original) & $q\ln q$ & $\{2,\,127,\,137,\,139,\,151,\,157\}$ \\
Stirling-truncated & $q\ln q - q + \tfrac{1}{2}\ln(2\pi q)$ & $\{2,\,389,\,397,\,401,\,409,\,421\}$ \\
Exact $\Gamma$     & $\ln\Gamma(q+1)$ & $\{2,\,389,\,397,\,401,\,409,\,421\}$ \\
\bottomrule
\end{tabular}
\end{table}

\begin{warnbox}
\textbf{Robustness finding}: The stable set is \textbf{not robust} under replacement of the
leading entropy by the exact Gamma entropy.  The set $\{2,127,137,139,151,157\}$ is
stable only under the leading $p\ln p$ approximation.  Under both Stirling-truncated and
exact $\Gamma$ entropies the stable set shifts entirely to $\{2,389,397,401,409,421\}$.

The physical reason: the Stirling correction $-q$ (representing the ``cost'' of
arranging $q$ quanta) shifts the effective stationarity from $p \approx e^5 \approx 148$
to $p \approx e^6 \approx 403$, relocating the stability window by one unit in $\ln p$.
\end{warnbox}

\subsection{Near-misses under exact $\Gamma$}

\begin{center}
\begin{tabular}{rrrrl}
\toprule
$p$ & $B(p)$ & $B_{\mathrm{low}}^\Gamma$ & $B_{\mathrm{high}}^\Gamma$ & Failure (margin) \\
\midrule
383 & 128.00 & 128.19 & 129.59 & lower ($-0.19$) \\
419 & 140.00 & 137.38 & 139.04 & upper ($-0.96$) \\
431 & 144.00 & 140.70 & 142.35 & upper ($-1.65$) \\
\bottomrule
\end{tabular}
\end{center}

\subsection{Verification: V-difference tables under exact $\Gamma$}

\subsubsection*{$p = 397$, $B = 132.667$ (exact $\Gamma$ model)}
\begin{center}
\begin{tabular}{rc}
\toprule
$q$ & $V_\Gamma(q;\,132.667) - V_\Gamma(397;\,132.667)$ \\
\midrule
383 & $+163.40$ \\
389 & $+53.53$ \\
\textbf{397} & \textbf{0} (minimum) \\
401 & $+13.20$ \\
409 & $+119.78$ \\
\bottomrule
\end{tabular}
\end{center}

\subsubsection*{$p = 421$, $B = 140.667$ (smallest margin, $\Delta_+ = 0.030$)}
\begin{center}
\begin{tabular}{rc}
\toprule
$q$ & $V_\Gamma(q;\,140.667) - V_\Gamma(421;\,140.667)$ \\
\midrule
409 & $+217.71$ \\
419 & $+19.66$ \\
\textbf{421} & \textbf{0} (minimum) \\
431 & $+1.80$ \quad \textbf{(nearest competitor)} \\
433 & $+22.23$ \\
\bottomrule
\end{tabular}
\end{center}

\subsection{Analytic criterion under exact $\Gamma$}

By the same argument as Section~\ref{sec:pss:s3}, the upper stability condition
under exact $\Gamma$ entropy gives to leading order:
\begin{equation}
  \frac{p+1}{3} \;<\; \frac{2p}{\ln p}
  \;\;\Longleftrightarrow\;\;
  (p+1)\ln p \;<\; 6p
  \;\;\Longrightarrow\;\;
  \ln p \;<\; 6, \quad p \;<\; e^6 \approx 403.
  \label{eq:pss:S8-upper}
\end{equation}

The analytic structure is identical to the original model but with the shift
$\ln p < 5 \to \ln p < 6$, reflecting the $-q$ Stirling correction absorbed into
the denominator of the stationarity value.

%──────────────────────────────────────────────────────────────────────────────
\section{Summary}
\label{sec:pss:summary}
%──────────────────────────────────────────────────────────────────────────────

\begin{resultbox}
\textbf{Main results}:
\begin{enumerate}
  \item \textbf{Characterisation}: A prime $p$ is prime-stable iff
        $B_{\mathrm{low}}(p) < (p+1)/3 < B_{\mathrm{high}}(p)$, where
        the bounds are set by the nearest primes $p_\pm$.
  \item \textbf{Analytic condition}: To leading order, stability requires
        $\ln p < 5$, i.e.\ $p < e^5 \approx 148$.
  \item \textbf{Complete stable set}: $\mathcal{S} = \{2, 127, 137, 139, 151, 157\}$,
        verified exhaustively up to $10^5$.
  \item \textbf{Finiteness}: The stable set is finite.  No prime beyond $157$ is
        prime-stable.  The required prime gap for stability grows as
        $g \sim p(\ln p - 5)/2$, far exceeding any plausible gap for $p \geq 300$.
  \item \textbf{No modular pattern}: The six stable primes lie in a 155-unit window.
        No congruence class or arithmetic structure selects them.
  \item \textbf{Modular reformulation}: With $B(p)=\mu(\Gamma_0(p))/3$, stability
        reads $\mu(\Gamma_0(p))(\ln p+1) < 6p$ — a condition on modular indices
        weighted by prime entropy.
  \item \textbf{Gamma-entropy reformulation} (Section~\ref{sec:pss:s8}): Replacing
        the leading entropy $q\ln q$ by the exact entropy $\ln\Gamma(q+1)$,
        or its Stirling truncation $q\ln q - q + \tfrac{1}{2}\ln(2\pi q)$,
        shifts the stability window from $\ln p < 5$ to $\ln p < 6$
        (i.e.\ $p \lesssim e^6 \approx 403$).
        The stable set under the exact model is $\{2,389,397,401,409,421\}$.
        The original set $\{2,127,137,139,151,157\}$ is \textbf{not robust}: it is
        an artefact of the leading-order entropy approximation.
\end{enumerate}
\end{resultbox}

\ifdefined\INCLUDEMODE
\else
  \end{document}
\fi
